%=======================================================================
%  CHAPTER 2 — PROBLEM SOLVING  (7 hrs)
%=======================================================================
\section{Problem Solving}

{\color{sub}\footnotesize 7 hrs \gap$\cdot$\gap asked in \mk{all 41 papers}
\gap$\cdot$\gap 4 syllabus sub-topics \gap$\cdot$\gap 366 marks, 10.8\,\% of every mark
ever set \gap$\cdot$\gap \mk{heavy solved-problem chapter}: see
\emph{Ch2 -- Problem Solving -- Solved Problems}.}

\lead{Read this first:}
\begin{itemize}
  \item \mk{One question dominates this chapter and the whole subject:}
        \textbf{solve a crypt-arithmetic puzzle by constraint satisfaction} \tS{31}.
        \textbf{31 of the 41 papers} set it, for 6 to 8 marks, and it is almost always
        preceded by a 1 or 2 mark \emph{``what is a CSP''} \tS{12}. Together that is
        \textbf{8 to 9 marks in nearly every paper}, from one technique.
  \item Thirteen different puzzles have been set. They are all worked in the companion
        file. \textbf{BASE $+$ BALL $=$ GAMES} (8 papers) and \textbf{SEND $+$ MORE $=$
        MONEY} (6) lead.
  \item \mk{Seven of the thirteen puzzles have more than one valid answer};
        ONE $+$ ONE $+$ TWO $=$ FOUR has \textbf{107}. So the marks are in
        \textbf{the column-by-column constraint reasoning}, not in the digits. Write the
        variables, the domains, the constraints and the carries, then reason. An answer
        with no working scores badly even when the arithmetic is right.
  \item The rest of the chapter is four small, reliable topics:
        \textbf{well-defined problem} \tF{7}, \textbf{state space and the steps of problem
        solving} \tF{5}, \textbf{production system} \tF{5}, and the
        \textbf{water-jug} formulation \tF{4}.
  \item Pairs the examiner reliably asks together: \textbf{well-defined problem $+$ a
        water-jug or crypt-arithmetic instance}, and \textbf{``what is a CSP'' $+$ the
        puzzle}. Almost nothing here is ever asked alone.
  \item \mk{Game playing is set up here and answered in Ch\,3.} This chapter
        owes you only the \emph{definition} of a game as a search problem; minimax and
        alpha-beta belong to \S3.4.
\end{itemize}

\hr

%-----------------------------------------------------------------------
\T{2.1 Problem solving and the state space}

\Q{Discuss the steps of problem solving. What is state space, successor function, goal test, path cost?}{\m{2}\m{3}}
{\color{sub}\footnotesize \tF{5} \yr{\textbf{\texttt{82 Bh}} \m{3}, \texttt{82 Ba} \m{2},
82 Ka \m{2}, \textbf{81 Ch} \m{3}, \textbf{\textit{76 Bh}} \m{3}} \gap$\cdot$\gap
\mk{a 2 or 3 mark opener}, never more. Four steps, five terms, done}

\sbs{%
\lead{The four steps of problem solving}
\begin{itemize}
  \item \textbf{1. Goal formulation} $-$ decide which world states count as success.
        Without it the agent has no way to prefer one action over another.
  \item \textbf{2. Problem formulation} $-$ decide what states and actions to consider:
        the initial state, the successor function, the goal test and the path cost.
  \item \textbf{3. Searching} $-$ find a sequence of actions leading from the initial
        state to a goal state, and choose the best such sequence. This is Ch\,3.
  \item \textbf{4. Execution} $-$ carry out the returned solution.
  \item \mk{Order matters and is often the mark}: goal before formulation,
        and the whole search happens \emph{before} any action is executed.
\end{itemize}
\lead{State space representation}
\begin{itemize}
  \item A \textbf{state space} is a set of \textbf{nodes}, one per state of the problem,
        joined by \textbf{arcs} representing the legal moves from one state to another,
        with a designated \textbf{initial state} and \textbf{goal state}.
  \item When a problem is defined with \emph{all} its states, it is a
        \textbf{complete state space}.
  \item A state space is drawn as a \textbf{directed graph}, or a \textbf{tree} when no
        state repeats. Nodes are search states, arcs are the operators.
\end{itemize}}{%
\lead{The five terms, with the 8-puzzle as the example}
\vspace{-1pt}
\begin{tabularx}{\linewidth}{@{}L{2.15cm}X@{}}
\toprule
\hrow \thead{Term} & \thead{Means, and for the 8-puzzle} \\
\midrule
\textbf{State} & A configuration of the world the agent may be in. \emph{The positions of the eight tiles and the blank.} \\
\textbf{Initial state} & Where the agent starts. \emph{The given starting board.} \\
\textbf{Successor function} & Given a state, returns the set of $\langle$action, resulting state$\rangle$ pairs reachable in one step. \emph{Move the blank left, right, up or down.} \\
\textbf{Goal test} & Decides whether a state is a goal. \emph{Does the board match the given goal board?} \\
\textbf{Path cost} & The sum of the step costs along a path. \emph{1 per move, so the cost is the number of moves.} \\
\bottomrule
\end{tabularx}
\vspace{2pt}
\begin{itemize}
  \item A \textbf{solution} is a sequence of actions from the initial state to a goal
        state. An \textbf{optimal solution} is one of lowest path cost.
\end{itemize}}

\penalty0\vspace{2pt}

\sbs{%
\lead{Why enumerating rules by hand does not work}
\begin{itemize}
  \item \textbf{The chess example} \yr{(\emph{Insights} \S2.1.1)}. A knight at
        $(4,4)$ can move to $(2,3)$, $(2,5)$, $(3,2)$, $(3,6)$, $(5,2)$, $(5,6)$, $(6,3)$
        and $(6,5)$.
  \item Writing that out as a rule for \emph{every} square gives an enormous rule set, and
        a chess position has around $10^{120}$ complete game paths.
  \item \mk{The fix, and the point of the example}: describe the move as a
        \textbf{pattern} rather than a list, i.e.\ current position $\pm$ (2 vertical $+$
        1 horizontal) or (1 vertical $+$ 2 horizontal). One rule replaces sixty-four.
  \item This is exactly why the next section exists: a \textbf{production system}
        is a set of patterns, not a table of cases.
\end{itemize}}{%
\lead{Problem space, and how to answer ``define a problem''}
\begin{itemize}
  \item \textbf{Problem space} $=$ every state reachable, plus every action that can be
        taken, for a given goal. \emph{Eg} the highways and cities of Nepal, when the
        problem is ``get from Kathmandu to Dharan''.
  \item \mk{The four-line answer that scores full marks} for
        \emph{``what are the criteria for defining a problem''} \yr{(82 Ka \m{2})}:
  \begin{itemize}
    \item define the \textbf{state space} the problem lives in,
    \item specify the \textbf{initial state},
    \item specify the \textbf{goal state} or a goal test,
    \item specify the \textbf{operators}, with their preconditions and effects.
  \end{itemize}
  \item Add \textbf{path cost} if the question mentions optimality or ``best''.
\end{itemize}}

%-----------------------------------------------------------------------
\T{2.2 Well-defined problems}

\Q{What makes a problem well defined? List the characteristics}{\m{2}\m{3}\m{4}\m{8}}
{\color{sub}\footnotesize \tF{7} \yr{\textbf{71 Bh} \m{3+4}, 75 Ba \m{2}, 80 Ash \m{2},
\textbf{78 Ch} \m{2}, \texttt{81 Ba} \m{2}, \textbf{75 Bh} \m{2}, 81 Ash \m{8}}
\gap$\cdot$\gap \mk{almost always a 2-mark opener} to a water-jug or
crypt-arithmetic question, so answer it in four bullets and move on. Only 81 Ash makes it
the whole 8-mark question}

\sbs{%
\lead{Definition}
\begin{itemize}
  \item A problem is \textbf{well defined} when the agent knows, before it starts,
        exactly what the problem consists of. Formally, it is composed of:
  \begin{itemize}
    \item an \textbf{initial state},
    \item the \textbf{actions}, given by a successor function,
    \item a \textbf{goal test}, and
    \item a \textbf{path cost} function.
  \end{itemize}
  \item Equivalently: the starting position, the allowable operations and the goal state
        are all clearly specified.
\end{itemize}
\lead{Characteristics \yr{(81 Ash asks for these by name, \m{8})}}
\begin{itemize}
  \item \textbf{Unambiguous}: one reading of what is being asked.
  \item \textbf{Complete}: everything needed to solve it is stated; nothing has to be
        assumed.
  \item \textbf{Verifiable}: the goal test can decide, mechanically, whether a state is a
        solution.
  \item \textbf{Finite and discrete} state space, so the states can be enumerated.
  \item \textbf{Deterministic operators} with known preconditions and effects.
  \item \textbf{Solvable in principle}: at least one action sequence reaches the goal.
\end{itemize}}{%
\lead{Well defined Vs ill defined, with examples}
\vspace{-1pt}
\begin{tabularx}{\linewidth}{@{}L{1.5cm}XX@{}}
\toprule
\hrow & \thead{Well defined} & \thead{Ill defined} \\
\midrule
Goal & Stated exactly & Vague or subjective \\
States & Enumerable & Open-ended \\
Operators & Fixed and known & Not fully known \\
Test & Mechanical & Needs judgement \\
Example & Water jug, 8-puzzle, chess, crypt-arithmetic & ``Write a good poem'', ``design a beautiful building'', medical diagnosis from an unclear history \\
\bottomrule
\end{tabularx}
\vspace{2pt}
\begin{itemize}
  \item \mk{The 5-litre / 2-litre trap} \yr{(81 Ash)}: the paper asks
        \emph{``is this a well-defined problem? Justify.''} The answer is \textbf{yes} $-$
        the jugs, the pump, the operators and the target are all given, so every clause
        above is satisfied. Say that, then solve it. The companion works the solution.
\end{itemize}}

%-----------------------------------------------------------------------
\T{2.3 Problem types and classification}

\Q{Problem types and how a problem is classified}{\m{2}\m{3}}
{\color{sub}\footnotesize
\pill{gdL}{gd}{syllabus, no PYQ yet} as a question of its own \gap$\cdot$\gap
but ``problem types'' is a printed syllabus line, and the \emph{ignorable / recoverable /
irrecoverable} split is worth knowing because it explains why the water jug is easy and
tic-tac-toe is not}

\sbs{%
\lead{By whether a move can be undone \yr{(\emph{Insights} \S2.2.1)}}
\begin{itemize}
  \item \textbf{Ignorable} $-$ intermediate steps can simply be ignored; a wrong move
        costs nothing. \emph{Eg} theorem proving, the \textbf{water jug}.
  \begin{itemize}
    \item Solvable with a simple control strategy, no backtracking needed.
  \end{itemize}
  \item \textbf{Recoverable} $-$ steps can be undone to get back to an earlier state.
        \emph{Eg} the \textbf{8-puzzle}.
  \begin{itemize}
    \item Needs backtracking, so the control strategy must keep a history.
  \end{itemize}
  \item \textbf{Irrecoverable} $-$ steps cannot be undone. \emph{Eg}
        \textbf{tic-tac-toe}, chess.
  \begin{itemize}
    \item Needs \textbf{planning}: consider the consequences before moving, because there
          is no second chance. This is why games need adversarial search (Ch\,3).
  \end{itemize}
  \item \textbf{Decomposable} $-$ the problem breaks into independent sub-problems whose
        answers combine. \emph{Eg} symbolic integration, word puzzles.
\end{itemize}}{%
\lead{By what the agent can see and predict}
\begin{itemize}
  \item \textbf{Single-state problem} $-$ deterministic and fully accessible. The agent
        always knows exactly which state it is in. \emph{Eg} chess with the board in view.
  \item \textbf{Multiple-state problem} $-$ deterministic but inaccessible. The agent
        knows the rules but not its exact state, so it reasons over a \emph{set} of
        possible states. \emph{Eg} walking across a dark room. It may have no sensors at
        all.
  \item \textbf{Contingency problem} $-$ non-deterministic and inaccessible. The
        environment is partially observable or the actions are uncertain, so the agent
        \textbf{must sense during execution} and cannot plan the whole path in advance.
  \item \textbf{Exploration problem} $-$ the state space itself is unknown. The agent
        discovers states and operators as it acts. \emph{Eg} a robot in an unmapped maze.
        \mk{The hardest of the four}, and the one that needs learning.
\end{itemize}}

%-----------------------------------------------------------------------
\T{2.4 Production systems}

\Q{Define a production system. What are its essential terms? Explain its role with an example}{\m{2}\m{4}\m{6}}
{\color{sub}\footnotesize \tF{5} \yr{70 Ma \m{2+6}, \textbf{\textit{71 Bh}} \m{4},
82 Ka \m{2}, \textbf{77 Ch} \m{2}, \textbf{75 Bh} \m{2+6}} \gap$\cdot$\gap
70 Ma pairs it with \emph{why is searching necessary in AI}; \textbf{75 Bh} pairs it with
\emph{well-defined problems}. It is the machinery behind the water-jug rules, so answer it
with those rules as the example}

\sbsr{0.54}{%
\lead{Definition and the four components}
\begin{itemize}
  \item A \textbf{production system} is a set of \textbf{rules}, each of the form
        $C_i \rightarrow A_i$:
  \begin{itemize}
    \item \textbf{LHS} $C_i$, the \textbf{pattern} or condition, which decides whether the
          rule applies;
    \item \textbf{RHS} $A_i$, the \textbf{action} to perform if it is applied.
  \end{itemize}
  \item The right-hand side is read as an \textbf{action recommendation}, not as a logical
        conclusion. That is what separates a production rule from an implication.
  \item \textbf{The four components}, and the answer to ``essential terms'':
  \begin{itemize}
    \item \textbf{Rule base} (production rules / rule memory) $-$ the rules themselves.
    \item \textbf{Working memory} (database) $-$ the current state of the world, updated
          as rules fire.
    \item \textbf{Control strategy} (inference engine) $-$ decides which applicable rule
          to fire when several match, and when to stop.
    \item \textbf{Rule applier} $-$ carries out the chosen action.
  \end{itemize}
  \item \textbf{The recognise-act cycle}: \textbf{match} the rules against working memory
        $\rightarrow$ \textbf{conflict resolution} picks one from the matching set
        $\rightarrow$ \textbf{act}, i.e.\ fire it and update working memory
        $\rightarrow$ repeat until the goal test passes or nothing matches.
\end{itemize}}{%
\lead{\mk{Two requirements of a good control strategy}}
\begin{itemize}
  \item \textbf{It must cause motion.} A strategy that can pick the same rule forever
        never leaves the initial state.
  \item \textbf{It must be systematic.} Choosing at random causes motion but wanders, and
        may revisit states endlessly.
  \item Both are one-line answers and both are asked.
\end{itemize}
\lead{Example: the vacuum world}
\vspace{-1pt}
\begin{tabularx}{\linewidth}{@{}L{2.4cm}X@{}}
\toprule
\hrow \thead{Pattern (LHS)} & \thead{Action (RHS)} \\
\midrule
$[A,\ \text{clean}]$ & move Right \\
$[A,\ \text{dirty}]$ & Suck \\
$[B,\ \text{clean}]$ & move Left \\
$[B,\ \text{dirty}]$ & Suck \\
\bottomrule
\end{tabularx}
\vspace{2pt}
\begin{itemize}
  \item \mk{The example that earns the marks} is the \textbf{water jug},
        whose eight production rules are in the companion. A question asking for a
        production system with an example almost always means that one.
  \item \textbf{Practical uses}: expert systems (Ch\,7), autonomous robots, and cognitive
        architectures modelling human reasoning such as ACT and SOAR.
\end{itemize}}

\penalty0\vspace{2pt}

\creamq{Problems solvable by production rules}{\m{6}}
{\color{sub}\footnotesize \yr{\textbf{75 Bh} \m{2+6}} \gap$\cdot$\gap
asked as ``explain about problems that can be solved using production rules with an
example''}

\sbs{%
\lead{Which problems suit a production system}
\begin{itemize}
  \item The state can be written down \textbf{compactly} $-$ the water jug's whole state
        is the pair $(x,y)$.
  \item The legal moves are \textbf{few, uniform and pattern-like}, so a handful of rules
        covers every case.
  \item Progress is made by \textbf{repeatedly rewriting the state}, not by evaluating a
        formula.
  \item The goal is a \textbf{recognisable pattern}, so the goal test is a match.
\end{itemize}}{%
\lead{The standard examples, and what each shows}
\begin{itemize}
  \item \textbf{Water jug} $-$ 8 rules, state $(x,y)$. The canonical answer.
  \item \textbf{8-puzzle} $-$ 4 rules (move the blank up, down, left, right), state $=$ the
        board. Shows a recoverable problem needing backtracking.
  \item \textbf{Missionaries and cannibals}, \textbf{farmer/wolf/goat/cabbage} $-$ rules
        with a \textbf{safety precondition}, which is what makes them interesting.
  \item \textbf{Tower of Hanoi} $-$ shows a decomposable problem.
  \item \textbf{Tic-tac-toe} $-$ shows an irrecoverable one, where rules alone are not
        enough and search over an opponent is needed.
\end{itemize}}

%-----------------------------------------------------------------------
\T{2.5 Constraint satisfaction problems}

\Q{What is a Constraint Satisfaction Problem? Discuss it and list its constraints}{\m{1}\m{2}\m{3}}
{\color{sub}\footnotesize \tS{12} \yr{\textbf{\texttt{81 Bh}}, \textbf{\texttt{80 Bh}},
\texttt{80 Ba}, 78 Ba, \textbf{80 Ch}, \textbf{\textit{74 Bh}} \m{2}, 79 Jth \m{1},
\textbf{79 Ch} \m{3}, \textit{74 Ma}, \textbf{\textit{72 Ash}}, 78 Po, \textbf{70 Bh}}
\gap$\cdot$\gap \mk{never asked alone}: it is the 1--3 mark preamble to a
crypt-arithmetic puzzle in every one of these papers, so keep it short and get to the
puzzle}

\sbsr{0.52}{%
\lead{Definition}
\begin{itemize}
  \item A \textbf{constraint satisfaction problem} is a search procedure that
        \textbf{operates in a space of constraints} rather than a space of states.
  \item States are defined by the \textbf{values assigned to variables}, and the goal test
        is a set of \textbf{constraints those values must obey}.
  \item Formally, a CSP is the triple:
  \begin{itemize}
    \item \textbf{Variables} $X = \{X_1, X_2, \dots, X_n\}$, the unknowns to be assigned.
    \item \textbf{Domains} $D = \{D_1, D_2, \dots, D_n\}$, the set of values each variable
          may take.
    \item \textbf{Constraints} $C = \{C_1, C_2, \dots, C_m\}$, the restrictions on which
          combinations of values are allowed.
  \end{itemize}
  \item \textbf{Solution}: an assignment of one domain value to every variable that
        violates no constraint. \mk{Write these three lines first in any
        crypt-arithmetic answer} $-$ they are worth marks on their own.
\end{itemize}
\lead{How it is solved: constraint propagation}
\begin{itemize}
  \item Constraints are \textbf{discovered and propagated as far as possible} through the
        system. Where propagation stalls, a \textbf{guess} is made and added as a new
        constraint, and propagation resumes.
  \item If the guess leads to a contradiction, \textbf{backtrack} and guess otherwise.
\end{itemize}}{%
\lead{\mk{The two termination conditions}}
\begin{itemize}
  \item \textbf{A contradiction is detected} $-$ no solution is consistent with the
        constraints now known.
  \item \textbf{Propagation has run out of steam} $-$ no further change can be made from
        current knowledge. Either a solution has been found, or a guess must be made.
  \item Both are asked as a pair. Two lines, two marks.
\end{itemize}
\lead{Types of constraint}
\begin{itemize}
  \item \textbf{Unary} $-$ restricts one variable. \emph{Eg} a leading letter is not 0.
  \item \textbf{Binary} $-$ relates two. \emph{Eg} two regions differ in colour.
  \item \textbf{Higher-order} $-$ relates three or more. \emph{Eg} a crypt-arithmetic
        column, $D + E = 10C_1 + Y$.
  \item \textbf{Global / AllDiff} $-$ all variables take different values. Every
        crypt-arithmetic puzzle carries this one.
\end{itemize}
\lead{Other standard CSP examples}
\begin{itemize}
  \item \textbf{Map colouring} $-$ no two adjacent regions share a colour.
  \item \textbf{$n$-queens} $-$ no queen attacks another.
  \item \textbf{Class scheduling}, \textbf{VLSI layout}, \textbf{Sudoku}.
\end{itemize}}

\penalty0\vspace{2pt}

\creamq{The crypt-arithmetic rules, and how to set the answer out}{\m{6}\m{7}\m{8}}
{\color{sub}\footnotesize \tS{31} \gap$\cdot$\gap
\mk{the highest-value single technique in the subject} \gap$\cdot$\gap
every puzzle is worked in \emph{Ch2 -- Problem Solving -- Solved Problems}}

\sbs{%
\lead{The rules, as every paper states them}
\begin{itemize}
  \item \textbf{Different letters denote different digits}, and identical letters denote
        the same digit. This is the AllDiff constraint.
  \item \textbf{No leading letter may be 0.}
  \item Digits come from $D = \{0,1,2,\dots,9\}$, so a puzzle with more than
        \textbf{10 distinct letters has no solution at all}.
  \item Any extra restriction the question prints overrides these.
\end{itemize}
\lead{\mk{The five-step method}} \yr{(\emph{Insights} \S2.3)}
\begin{itemize}
  \item \textbf{1.} Write the sum in columns with a \textbf{carry variable}
        $C_1, C_2, \dots$ above each column boundary.
  \item \textbf{2.} List $X$, the set of variables.
  \item \textbf{3.} List $D$, the domain.
  \item \textbf{4.} List $C$: the leading-letter constraints, the AllDiff constraint, and
        \textbf{one equation per column},
        $\text{(column sum)} + C_{\text{in}} = 10\,C_{\text{out}} + \text{(result digit)}$.
  \item \textbf{5.} Reason column by column, redrawing the sum with the digits filled in
        after each deduction, until every letter is fixed. Close with the arithmetic
        check.
\end{itemize}}{%
\lead{The three deductions that crack most puzzles}
\begin{itemize}
  \item \textbf{A carry out of the leftmost column can only be 1.} If the answer is one
        digit longer than the addends, that leading digit \textbf{is} 1. This is what
        gives $M = 1$ in SEND $+$ MORE $=$ MONEY immediately.
  \item \textbf{A letter added to itself.} $2N$ ends in the same letter only when that
        letter is 0 or 5, which is how TEN $+$ TEN $+$ FORTY $=$ SIXTY is cracked.
  \item \textbf{A column where a letter repeats.} If $E + O + C_2 = N$ and also
        $N + R + C_1 = E + 10$, substitute one into the other; the unknowns cancel.
\end{itemize}
\lead{\mk{Two things that lose marks}}
\begin{itemize}
  \item \textbf{Giving only the digits.} Most of these puzzles have many solutions, so an
        unjustified assignment cannot be marked. Show the columns.
  \item \textbf{Forgetting the AllDiff check at the end.} It is the single commonest way a
        nearly-right answer fails, and it is what rules out $B = 9$ in the
        $AB + CD = AAA$ question \yr{(\textbf{\textit{76 Bh}})}.
\end{itemize}}

%-----------------------------------------------------------------------
\T{2.6 Game playing}

\Q{Formulate a game as a search problem}{\m{2}\m{7}}
{\color{sub}\footnotesize \tP{1} \yr{\textit{72 Ma} \m{7}, tic-tac-toe state space}
\gap$\cdot$\gap \mk{the algorithms live in \S3.4}, not here. This section only
owes you the formulation; minimax and alpha-beta are Ch\,3, where they are asked 14 times}

\sbs{%
\lead{A game, defined as a search problem}
\begin{itemize}
  \item \textbf{Initial state} $-$ the board as set up, plus whose turn it is.
  \item \textbf{Operators} $-$ the legal moves.
  \item \textbf{Terminal test} $-$ says when the game is over.
  \item \textbf{Utility / payoff function} $-$ a number on each terminal state saying who
        won and by how much. \emph{Eg} $+1$ win, $0$ draw, $-1$ loss.
  \item \mk{What makes it different from ordinary search}: there is an
        \textbf{opponent}, choosing to make things as bad as possible for you. So a plan
        cannot be a fixed sequence of moves; it must be a \textbf{strategy}, an answer for
        every reply.
  \item Two-player, perfect information, zero-sum games are solved by
        \textbf{minimax} \yr{(\S3.4)}, made efficient by \textbf{alpha-beta pruning}.
\end{itemize}}{%
\lead{Why games were chosen as an AI testbed}
\begin{itemize}
  \item Rules are \textbf{simple and unambiguous}, so no knowledge acquisition problem
        gets in the way.
  \item The state is \textbf{fully observable}, so no perception problem either.
  \item Yet the \textbf{search space is enormous}, so brute force fails and real technique
        is needed. That combination is rare and useful.
  \item Performance is \textbf{measurable} against human players.
\end{itemize}
\lead{Production systems Vs game playing}
\begin{itemize}
  \item A production system fires rules against \textbf{its own} working memory.
  \item A game must additionally model \textbf{another agent choosing against it}, which
        no rule set alone can express. That is precisely why Ch\,3 needs a separate family
        of algorithms.
\end{itemize}}
